%% Paper 1: sensitivity of a gridded European heat-wave catalogue to its
%% detection thresholds.
%%
%% Needs the AMS LaTeX package version 6. Two files are NOT in this repository
%% and have to be downloaded from the AMS LaTeX submissions page, which refuses
%% automated requests, then placed beside this file:
%%
%%     ametsocV6.cls    the document class
%%     ametsocV6.bst    the bibliography style
%%
%% https://www.ametsoc.org/ams/publications/author-information/latex-author-info/
%%
%% Editorial Manager builds the PDF from the source, so do not upload a PDF of
%% your own. Figure files it receives must carry flat names with no directory
%% path, no underscores and no periods before the extension, which is why the
%% figure macros below expect fig1.pdf and not plots/detection/europe_....pdf.
%%
%% Compile: pdflatex, bibtex, pdflatex, pdflatex.

%% PREVIEW: swap back to ametsocV6 once the AMS files are here.
\documentclass{amsdraft}

\usepackage{amsmath,amsfonts,amssymb,bm}
\usepackage{mathptmx}
\usepackage{newtxtext}
\usepackage{newtxmath}
\usepackage{graphicx}

\title{TITLE DEFERRED: written last, from the finished results}

%% AMS designates ONE corresponding author on the title page. The
%% corresponding author in Editorial Manager can be a different person.
\authors{Isaac Adoboe\correspondingauthor{Isaac Adoboe, ieadoboe@mun.ca}
and Elsa Cardoso-Bihlo}

\affiliation{Memorial University of Newfoundland, St. John's,
Newfoundland and Labrador, Canada}

\abstract{ABSTRACT DEFERRED: written last, once the results and conclusions
are settled.}

\begin{document}

\maketitle

%% Optional for a journal article, required for BAMS, at most 120 words.
% \statement
% Significance statement here.

\section{Introduction}
\label{sec:intro}

%% DRAFT, written before the results. Revisit the last paragraph once
%% section 3 exists, so the question it poses matches what was answered.

European summer heat waves are among the most damaging weather extremes the
continent experiences. The summer of 2003 was statistically extremely
unlikely even when the observed warming is taken into account
\citep{Schar2004}, and more than 70\,000 additional deaths occurred during it
\citep{Robine2008}; the summer of 2010, centred on eastern Europe and Russia,
then exceeded 2003 in both amplitude and spatial extent
\citep{Barriopedro2011}; and the summer of 2022 was associated with more than
60\,000 heat-related deaths \citep{Ballester2023}. Over the past four decades,
heat-wave trends over Europe have been three to four times faster than over
the rest of the northern midlatitudes \citep{Rousi2022}.

Deciding which of these events were the most severe is usually the job of a
catalogue. Ranked lists of the most severe European heat waves
\citep{Russo2015,Becker2022} and classifications of their spatial patterns
\citep{Stefanon2012} determine which events are studied in depth.
Extreme-event attribution starts from the choice of an event and its
definition \citep{Philip2020}, and physical storylines of extreme heat are
constructed around a chosen case \citep{WhittakerDiLuca2026}. In each of these
uses the catalogue is an input rather than a result: once an event has been
selected as the most severe of its season, the work that follows inherits
that selection without revisiting it.

There is, however, no single definition of a heat wave
\citep{PerkinsAlexander2013,Perkins2015}. Indices developed for station
records count hot days at a point \citep{Zhang2011}, whereas a gridded
catalogue must also decide how large a hot region has to be before it
constitutes an event \citep{Stefanon2012,Lhotka2015}. Every such detector has
parameters with no physically correct value, among them the percentile that
defines a hot day and the minimum area that defines a coherent event.
Comparisons between definitions show that the frequency, trends and
geographic patterns of heat waves depend strongly on the index adopted
\citep{Smith2013}, and that the ranking of historical European events shifts
between indices \citep{Becker2022}. Within a single event-based detector,
the area threshold, connectivity and overlap criterion have been shown to
alter heat-wave metrics and their trends \citep{Zhang2025}. Those studies
concern aggregate statistics. Less attention has been paid to a narrower
question that matters directly to anyone selecting cases: for a single
detector, how much does the ranking of its individual events depend on its
own free parameters? Choosing the top event of a catalogue implicitly assumes
that the answer is very little.

This paper tests that assumption. We construct a gridded CTX90pct catalogue
of European heat waves from ERA5 for the 2026 season and for two contrasting
seasons, 2003 and 2022, and repeat the detection for six combinations of
percentile and minimum area. We then ask whether the most severe events of
each season keep their rank across configurations and, where they do not,
which step of the detector is responsible. Section~\ref{sec:methods} describes
the data, the detector and the sensitivity experiment;
section~\ref{sec:results} presents the catalogues and their response to the
parameters; section~\ref{sec:discussion} considers what that response implies
for case selection; and section~\ref{sec:conclusions} summarizes.

\section{Data and methods}
\label{sec:methods}

\subsection{Reanalysis data and domain}
\label{sec:data}

We built the catalogue from daily maximum 2-m temperature in the ERA5
reanalysis \citep{Hersbach2020}, retrieved from the Copernicus Climate Data
Store as the derived daily-statistics product, which aggregates the hourly
reanalysis to a daily maximum before subsetting in space. The domain spans
35--72$^\circ$N and 25$^\circ$W--45$^\circ$E on the native 0.25$^\circ$ grid,
which is 149 latitudes by 281 longitudes, and the calendar window runs from 1
May to 30 September, giving 153 daily fields per season.

Two further seasons, 2003 and 2022, were retrieved on the same grid over the
same calendar window and processed with an identical configuration. They were
chosen because their heat was organized differently from one another and from
2026: the severity of 2003 is dominated by a single August event over France,
whereas 2022 combines a sequence of western European events with a separate and
longer-lived northeastern one. Testing the detector against seasons of
different shape is what distinguishes a property of the method from a property
of one summer.

Thresholds were computed from a 1991--2020 baseline, the current World
Meteorological Organization reference normal. Ocean cells were removed with a
land mask taken from the 0.1$^\circ$ K\"oppen-Geiger map of \citet{Beck2023},
in which class zero denotes ocean, sampled to the ERA5 grid by nearest
neighbour. Deriving the mask from that map rather than from a separate
land-sea product keeps the definition of land consistent with the climate
classification applied to each event in section~\ref{sec:events}.

\subsection{Heat-wave index}
\label{sec:index}

The index is CTX90pct as defined by \citet{PerkinsAlexander2013}, applied to
every land cell independently. For each cell and each calendar day $d$ we
computed the 90th percentile of daily maximum temperature over the baseline,
pooling every baseline day whose day-of-year lies within $\pm$7 days of the
target, with distances in day-of-year evaluated circularly. Each threshold
$T_{90}(d,i,j)$ is therefore estimated from 30 years of 15 days, or 450
samples, which is what makes a 90th percentile stable enough to threshold
against, and it varies smoothly through the season and with position rather
than imposing one fixed temperature across a domain that spans 37 degrees of
latitude.

A cell-day was marked hot when its daily maximum exceeded its own threshold,
and we retained the exceedance anomaly
$T_{\max}(t,i,j) - T_{90}(d_t,i,j)$ for every such cell-day. A hot cell-day was
kept only if it belonged to a run of at least three consecutive hot days at
that cell. We implemented the persistence requirement as a binary opening
along the time axis with a structuring element of length three: the erosion
removes every day belonging to a run shorter than three days, and the dilation
restores in full the runs that survive.

\subsection{Spatial coherence}
\label{sec:coherence}

Applying CTX90pct to a grid yields a boolean field of cell-days, which is not
yet a catalogue of events. Turning it into one requires grouping cell-days
that belong together, and the order of operations at this stage determines
whether that grouping is meaningful.

A 90th-percentile threshold places approximately one tenth of all cell-days
above it by construction, so a persistent background of scattered exceedances
covers the domain on any given day. Labelling that field directly allows
isolated cells to act as bridges between unrelated patches and to chain them
into a single object: in our own testing, a planted five-day event was
recovered as a twelve-day one. We therefore filtered each day for spatial
extent before any labelling in time. Within each day, hot cells were grouped
into connected patches under eight-connectivity, the area of each patch was
accumulated from the spherical cell areas
\begin{equation}
A(i) = R^2 \, \Delta\phi \, \Delta\lambda \, \cos\phi_i ,
\label{eq:area}
\end{equation}
with $R = 6371.0088$ km and $\Delta\phi$, $\Delta\lambda$ the grid spacing in
radians, and every patch whose area fell below a minimum $A_{\min}$ was
discarded. Weighting by area rather than counting cells matters across this
domain, where a cell at 72$^\circ$N covers about 38\% of the area of one at
35$^\circ$N. The reference configuration uses $A_{\min} = 50\,000$ km$^2$,
comparable to the area of Slovakia. Section~\ref{sec:sensitivity} treats that
value, and the percentile, as free parameters rather than as settled choices.

\subsection{Event identification and attributes}
\label{sec:events}

Events are the connected components of the surviving cell-days in the
three-dimensional field of time, latitude and longitude, labelled with the
\texttt{scipy.ndimage} implementation \citep{Virtanen2020}. The connectivity
structure joins each cell to its eight spatial neighbours within the same day
and to itself on the preceding and following day, but not to cells that are
diagonal in space and time at once. Full 26-neighbour connectivity merges
events that only brush past one another. A labelled component was retained as
an event when its largest daily area reached $A_{\min}$, which removes the
single-cell specks that otherwise accumulate along coastlines and over
mountains.

Each event carries a severity,
\begin{equation}
S = \sum_{(t,i,j)\,\in\,E}
    \left[ T_{\max}(t,i,j) - T_{90}(d_t,i,j) \right] A(i) ,
\label{eq:severity}
\end{equation}
the exceedance anomaly summed over every cell-day in the event $E$ and
weighted by cell area, in units of $^\circ$C km$^2$ day. Severity is the
quantity by which the catalogue is ranked throughout this paper. It is a
gridded analogue of the heat wave magnitude index of \citet{Russo2014}: an
event scores highly by being hot, large and long together, so that a brief
record at a single location does not outrank a continental month.

For each event we also recorded the start, end and duration; the peak and mean
daily areas; the highest daily maximum temperature and the largest exceedance
reached anywhere within it; and the peak day, defined as the day on which the
mean exceedance across the event's cells was greatest, which is the day it was
most intense rather than the day it was most widespread. The centroid is the
area-weighted mean position of the cells belonging to the event on its peak
day. Each event was finally assigned the modal K\"oppen-Geiger class of the
land cells it covered at any point in its life, again from \citet{Beck2023}.

\subsection{Sensitivity experiment}
\label{sec:sensitivity}

Two parameters of the detector have no correct value. The percentile sets how
far above its own climatology a cell must be to count as hot, and the minimum
area $A_{\min}$ sets how large a patch must be on a given day to count at all.
Both are conventions, and published catalogues adopt different ones. We
therefore repeated the entire detection for every combination of three
minimum areas (25\,000, 50\,000 and 100\,000 km$^2$) and two percentiles (the
90th and the 95th), giving six configurations, and applied all six to each of
the three seasons. The configuration with $A_{\min} = 50\,000$ km$^2$ and the
90th percentile is the reference, and the remaining five are compared against
it.

Comparison requires matching events between runs, because the labelling
numbers components in scan order and no identifier survives a change of
configuration. For each of the six most severe reference events we searched
the target catalogue for a counterpart overlapping it in time by at least one
day and lying within 10$^\circ$ of the reference centroid, with longitude
differences scaled by the cosine of the reference latitude. Among the
candidates we took the one sharing the most days with the reference, breaking
ties by centroid distance. The quantity we report is the change in severity
rank, defined as the reference rank minus the rank the matched event holds in
the target catalogue, so that a positive value means the event moved up the
ranking.

One property of this procedure deserves note. It operates on dates and
centroids rather than on the label fields themselves, so where a configuration
splits a reference event into two, only the larger fragment is reported. The
converse case, in which a configuration detects two reference events as a
single object, appears as two reference events matching the same target event
and is identified as such in section~\ref{sec:results}.

Finally, 2003 lies inside the 1991--2020 baseline and so contributes to the
thresholds against which it is tested, which raises its own thresholds and
makes its catalogue conservative relative to those of 2022 and 2026. We
retain the common baseline so that all three seasons are measured against one
climatology, and return to the consequences in section~\ref{sec:discussion}.

\section{Results}
\label{sec:results}

%% Numbers come from data/heatwave_events_europe_<year>.csv and
%% data/catalogue_sensitivity_europe_<year>.csv (2026 rerun 2026-10-07 on
%% the full season to 30 September). Figure sources, renamed flat for AMS:
%%   fig1.pdf  plots/detection/europe_2026_fig1b_map_borders.pdf
%%   fig2.pdf  plots/detection/europe_2026_fig2_timeline.pdf
%%   fig3.pdf  plots/detection/europe_2026_fig6_footprints.pdf
%%   fig4.pdf  plots/detection/europe_sensitivity_rank_changes.pdf

\subsection{The 2026 catalogue}
\label{sec:cat2026}

The reference configuration detected 39 events between 1 May and 30
September 2026 (Figs.~\ref{fig:map} and~\ref{fig:timeline}). Their
severities span almost three orders of magnitude, from
$1.5\times10^{5}$ to $1.2\times10^{8}$\,$^\circ$C\,km$^2$\,day, and the
total is concentrated in a few of them. The most severe event alone
accounts for 25\% of the summed severity of the season, the three most
severe for 51\% and the six most severe for 75\%; seven events are enough
to reach 80\%. The remaining 32 are short and small: the median duration
across the catalogue is six days.

The most severe event ran from 12 June to 3 July. It covered 2.98 million
km$^2$ on its largest day, the greatest extent of any event in the season,
spread from the Bay of Biscay to Poland (Fig.~\ref{fig:footprints}), and
reached the largest exceedance in the catalogue, 14.1$^\circ$C above the
local threshold. Its hottest cell, at 43.6$^\circ$C, lay in southwestern
France, more than 1\,000 km west of its centroid in southern Germany. The
second event, 27 July to 19 August, was centred on southern France,
with its longest-lived cells over the Alps and northern Italy, and peaked
at 1.96 million km$^2$; the third, 20 May to 2 June,
covered Iberia, France and central Europe; and the fourth, 30 June to
24 July, was confined largely to France, Iberia and western
Italy. Figure~\ref{fig:footprints} shows that France lies
inside the footprints of all four, and that the first and fourth overlap
in time as well as in space: the fourth began on 30 June, four days before
the first ended. The season's severity was therefore delivered as a
sequence of waves over largely the same ground rather than as separate
regional episodes, a property that section~\ref{sec:sens} shows to matter.

Counted by events, the catalogue is dominated by the temperate oceanic
class Cfb (12 of 39 events), but weighted by severity it is dominated by
the humid continental class Dfb, whose eight events carry 50\% of the
total against 35\% for Cfb. Absolute and relative heat also rank the
events differently. The highest daily maximum in the catalogue,
49.2$^\circ$C, was recorded in an event at the southeastern corner
of the domain whose largest exceedance was only 4.1$^\circ$C, whereas the
largest exceedance belongs to the central European event of June. Across
all 39 events the Spearman rank correlation between peak temperature and
peak exceedance is 0.11.

\subsection{The 2003 and 2022 catalogues}
\label{sec:cat_other}

The same configuration detected 41 events in 2003 and 42 in 2022, a count
close to the 39 of 2026. Severity is similarly concentrated: the three
most severe events hold 55\% of the season total in 2003 and 62\% in
2022. The seasons differ in where that severity sat. In 2003 the most
severe event was the August heat wave over France, 29 July to 15 August,
which accounts for 23\% of the season and reached 46.9$^\circ$C. The
second, centred on the Pannonian basin, begins on 1 May, the first day of
the analysis window, and is therefore truncated. In 2022 the two most
severe events were both northeastern. The first, centred on northwestern
Russia, persisted for 39 days from 29 July to 5 September and holds 31\%
of the season total; the second, 21 June to 11 July, was centred on the
Baltic. The western European heat waves of that summer rank third and
fourth, with peak temperatures of 44.5 and 43.7$^\circ$C; the third
reached the largest exceedance of the season, 15.0$^\circ$C. In neither
season do the leading events share a region in consecutive weeks as the
2026 events do.

\subsection{Sensitivity to percentile and minimum area}
\label{sec:sens}

The number of events responds strongly to both parameters
(Fig.~\ref{fig:sensitivity}). At the 90th percentile, halving $A_{\min}$
from 50\,000 to 25\,000 km$^2$ raised the count by 61\% in 2003, 74\% in
2022 and 72\% in 2026, and doubling it to 100\,000 km$^2$ removed between
24\% and 41\%. Raising the percentile to the 95th removed further events
at every area, and cut the severity of each season's leading event to
between a third and three fifths of its reference value, which is
expected because both the set of hot cell-days and the
exceedance of each one shrink. None of this is a defect in itself: a
catalogue ranked by severity is unaffected by small events entering or
leaving it, provided the leading events keep their order.

Every one of the six leading reference events of every season had a
counterpart in all six configurations, so no leading event disappeared.
Their order is a different matter. In 2003 and 2022 the three most severe
events held their ranks in all six configurations, and nothing moved at
all at the 90th percentile. All the reordering in those two seasons
occurred at the 95th percentile and below the third rank. In 2003 the
June event fell from fourth to sixth, and the two Fennoscandian events of July
each rose one place. In 2022 the late-May event over Anatolia fell from sixth to
ninth, or to tenth at 100\,000 km$^2$. These shifts are unchanged by the
minimum area, which shows they come from the percentile alone: the events
concerned were modest in exceedance and lost a larger share of their
severity than their neighbours when the threshold rose.

The 2026 ranking responded to both parameters and reached its second rank.
At the 90th percentile and 25\,000 km$^2$ the June event and the event of
30 June to 24 July were detected as a single object running 44 days from
12 June to 25 July, whose severity,
$1.67\times10^{8}$\,$^\circ$C\,km$^2$\,day, is the sum of the two. That
object still ranks first, but the fourth reference event no longer exists
as a separate event, and the fifth and sixth each rise one place into the
gap. At the larger minimum areas the daily patches that connect the two
waves during their four days of overlap fall below $A_{\min}$, and the
two are separated again. At the 95th percentile, the second and third
events exchanged ranks at 25\,000 and 50\,000 km$^2$: their severities
differ there by less than 4\%. At 100\,000 km$^2$ the second event broke
apart. Its largest fragment, 8 to 16 August, ranks only seventh, a fall of
five places, and the reference ranks four to six rearrange around the
vacancy. Because the matching reports only the larger fragment of a
split event (section~\ref{sec:sensitivity}), the remaining fragments
appear elsewhere in that catalogue as separate, smaller events.

Across the three seasons the leading event was never displaced. Below it,
2003 and 2022 kept their top three under every configuration, whereas 2026
changed its second-ranked event under three of the five alternatives and
changed the composition of its leading events under a fourth. The seasons
whose severity was concentrated in events that were separated in space or
time produced stable rankings; the season whose leading events followed
one another over the same region did not.

\begin{figure}[tbp]
  \centering
  \includegraphics[width=\textwidth,height=0.75\textheight,keepaspectratio]{fig1}
  \caption{Number of days each land cell spent inside any of the 39 events
    detected in 2026 with the reference configuration (90th percentile,
    $A_{\min} = 50\,000$ km$^2$). Numbers mark event centroids and give
    the severity rank.}
  \label{fig:map}
\end{figure}

\begin{figure}[tbp]
  \centering
  \includegraphics[width=\textwidth,height=0.75\textheight,keepaspectratio]{fig2}
  \caption{Duration of each 2026 event, ordered by start date and shaded by
    severity. The tick inside each bar marks the peak day. Labels give the
    severity rank and the modal K\"oppen-Geiger class.}
  \label{fig:timeline}
\end{figure}

\begin{figure}[tbp]
  \centering
  \includegraphics[width=\textwidth,height=0.75\textheight,keepaspectratio]{fig3}
  \caption{Footprints of the eight most severe events of 2026: the number
    of days each cell spent inside the event. The star marks the cell that
    recorded the event's highest daily maximum temperature.}
  \label{fig:footprints}
\end{figure}

\begin{figure}[tbp]
  \centering
  \includegraphics[width=\textwidth,height=0.75\textheight,keepaspectratio]{fig4}
  \caption{Change in severity rank of the six most severe reference events
    of each season under each configuration, defined as reference rank
    minus new rank, so that a positive value is a rise. Blank cells are
    unchanged. Column heads give $A_{\min}$, and the numbers beneath each
    panel give the total number of events detected. The blue outline joins
    reference events that a configuration detects as a single event.}
  \label{fig:sensitivity}
\end{figure}

\clearpage

\section{Discussion}
\label{sec:discussion}

%% TODO. Written third.

\section{Conclusions}
\label{sec:conclusions}

%% TODO.

\datastatement
%% TODO. AMS requires this section. It needs the Zenodo DOI for the deposited
%% catalogues, the ERA5 CDS citation, and the code repository.

%% PREVIEW: swap back to ametsocV6 once the AMS files are here.
\bibliographystyle{plainnat}
\bibliography{references}

\end{document}
